\documentclass[12pt,a4paper]{article}
\usepackage[utf8]{inputenc}
\usepackage[portuguese]{babel}
\usepackage{amsmath}
\usepackage{graphicx}
\usepackage{booktabs}

\title{Análise Analítica dos Pré-Laboratórios \\ \large Diodos e Retificadores}
\author{Cleber da Silva Alves}
\date{\today}

\begin{document}

\maketitle

\section{Pré-Laboratório 4: Circuitos Básicos}

\subsection{Modelagem Simplificada ($V_{DO} = 0,7$ V)}
Para os circuitos retificadores, detectores de pico e grampeadores, o diodo é modelado como uma chave em série com uma fonte de $0,7$ V.

\begin{itemize}
    \item \textbf{Retificador de Meia Onda:} $V_o = V_i - 0,7$ para $V_i > 0,7$V; caso contrário $V_o = 0$.
    \item \textbf{Detector de Pico:} O capacitor carrega até $V_p - 0,7$V e mantém o valor (idealmente).
    \item \textbf{Grampeador Positivo:} Adiciona um nível DC ao sinal de modo que o pico negativo fique em $-0,7$V. $V_o = V_i + (V_p - 0,7)$.
\end{itemize}

\subsection{Determinação de $\eta$ e $I_S$}
Dados: $V_T = 25$ mV. Usando dois pontos da tabela ($I_{D1}, V_{D1}$) e ($I_{D2}, V_{D2}$):
\begin{equation}
\eta = \frac{V_{D2} - V_{D1}}{V_T \ln(I_{D2}/I_{D1})}
\end{equation}
Para $I_D = 0,1$ mA ($V_D = 500$ mV) e $I_D = 1,0$ mA ($V_D = 601$ mV):
\begin{equation}
\eta = \frac{601 - 500}{25 \ln(10)} \approx \frac{101}{57,56} \approx 1,75
\end{equation}
Extrapolando para $I_S$: $I_S = I_D e^{-V_D/(\eta V_T)} \approx 0,1 \text{mA} \cdot e^{-500/(1,75 \cdot 25)} \approx 1,05 \text{ nA}$.

\section{Pré-Laboratório 5: Retificador de Baixa Tensão}

\subsection{Dedução da Equação de Tensão de Saída}
A equação (2) do roteiro descreve a tensão de saída $V_O$ em função da corrente de carga $I_L$ e das tensões de entrada $V_H$ e $V_L$:
\begin{equation}
V_O = \frac{V_H + V_L}{2} + \eta V_T \ln \left[ \frac{\cosh(\frac{V_H - V_L}{2 \eta V_T})}{1 + \frac{I_L}{I_S}} \right]
\end{equation}
Esta equação é fundamental quando a amplitude do sinal é comparável à queda de tensão no diodo, onde o modelo de $0,7$V falha.

\section{Pré-Laboratório 6: Regulador Zener}

\subsection{Projeto do Resistor de Proteção $R_P$}
Dados: $V_C = 8,5$ V, $V_Z = 6,2$ V, $I_{Zmax} = 7$ mA.
\begin{equation}
R_P = \frac{V_C - V_Z}{I_{Zmax}} = \frac{8,5 - 6,2}{0,007} \approx 328,5 \Omega
\end{equation}
Valor comercial E-12 superior: \textbf{$330 \Omega$}.
Potência no Zener: $P_Z = V_Z \cdot I_{Zmax} = 6,2 \cdot 7\text{m} = 43,4$ mW (Bem abaixo dos 500 mW típicos).

\end{document}
